% THIS DOCUMENT WAS COMPUTER GENERATED FROM  diploma-prazno
% IT REQUIRES TO BE RUN FROM TERMINAL SINCE GLOBAL .vsc SETTINGS

\documentclass{beamer}

\usetheme{Warsaw}

\setbeamertemplate{itemize items}[circle]


\usepackage[slovene]{babel}
\usepackage[utf8]{inputenc}
\usepackage[T1]{fontenc}
\usepackage{lmodern}

\usepackage{amsfonts, amsmath, amsthm}
\usepackage{enumitem}

% --- Teorematična okolja ---
\newtheorem{izrek}{Izrek}
\newtheorem{trditev}[izrek]{Trditev}
\newtheorem{posledica}[izrek]{Posledica}
\newtheorem{definicija}[izrek]{Definicija}

% --- Ukazi ---
\newcommand{\red}{\operatorname{red}}
\newcommand{\srg}{\operatorname{srg}}
\newcommand{\diam}{\operatorname{diam}}
\newcommand{\N}{\mathbb{N}}
\newcommand{\ZZ}{\mathbb{Z}}
\newcommand{\reach}{\operatorname{Reach}}
\newcommand{\gr}{\mathscr{G}}

\newcommand{\A}{\mathcal{A}}
\newcommand{\II}{\mathbf{I}}
\newcommand{\IIw}[2]{#1 \;\II\; #2}
\newcommand{\IIwd}[2]{#1 \;\II^D\; #2}
\newcommand{\PP}{\mathcal{P}}
\newcommand{\LL}{\mathcal{L}}

% --- Naslov ---
\title{Mooreovi grafi in posplošeni večkotniki}
\author{Luka Lavš}
\date{\today}

\begin{document}

% --- Title ---
\begin{frame}
    \titlepage
\end{frame}

% ------------------------------
\section{Mooreovi grafi}

\begin{frame}
    \frametitle{Uvodne definicije}
    \begin{itemize}
        \item[\textbullet] Graf $G$ ima \textbf{premer} $D$, če je največja razdalja med  
        katerima koli dvema vozliščema enaka $D$.
        \item[\textbullet] Graf $G$ ima \textbf{maksimalno stopnjo} $\Delta$, če je
        $\Delta = \max_{v \in V(G)} \deg(v)$.
        \item[\textbullet] Graf $G$ ima \textbf{ožino} $g$, če je dolžina 
        najkrajšega cikla v $G$ enaka $g$.
        \item[\textbullet] Graf $G$ je \textbf{$\Delta$-regularen}, če ima vsako 
        vozlišče stopnjo $\Delta$.
    \end{itemize}
\end{frame}

\begin{frame}
\frametitle{Definicija Mooreovega grafa}
\begin{definicija}
Naj bo $G=(V,E)$ graf s premerom $D$ in maksimalno stopnjo $\Delta$.  
Če velja
\[
|V| = 1 + \Delta + \Delta(\Delta-1) + \dots + \Delta(\Delta-1)^{D-1},
\]
pravimo, da je $G$ \emph{Mooreov graf}.
\end{definicija}
\end{frame}

\begin{frame}
\frametitle{Lastnosti Mooreovih grafov}
\begin{trditev}
Moorejev graf $G$ je $\Delta$-regularen. Med vsakima dvema vozliščema obstaja natanko ena pot dolžine največ $D$.
\end{trditev}

\begin{posledica}
Ožina Moorejevega grafa je $g = 2D + 1$.
\end{posledica}
\end{frame}


\begin{frame}
    \frametitle{Matrična enačba Mooreovih grafov}
    \begin{trditev}
    Za regularen Moorejev graf $G$ stopnje $\Delta$ z matriko s
sosednosti $A$ velja:
\[G_{\Delta, D} = J,\]    
kjer je $J$ matrika samih enic in 
$G_{\Delta, D} = A G_{\Delta, D-1} - 
(\Delta - 1) G_{\Delta, D-2}$, 
$G_{\Delta, 1} = A + I$ in $G_{\Delta, 0} = I$.
    \end{trditev}
\end{frame}

\begin{frame}
    \frametitle{Redkost in razširitve Mooreovih grafov}
    \begin{itemize}
        \item[\textbullet] Iskanje grafov, ki dosežejo  
        red $M_{\Delta, D} - \delta$ za "majhen"~$\delta$. 
        \item[\textbullet] Spodnja meja za red grafov ožine $g=2D+1$
        ter maksimalne stopnje $\Delta$, razširja zgornjo mejo 
        za Mooreove grafe. 

    \end{itemize}
            \begin{trditev}
            \begin{equation*}
        |V(G)| \ge
        \begin{cases}
        \frac{\Delta(\Delta-1)^{(g-1)/2}-1}{\Delta-2}, & g \text{ liho},\\[1em]
        2\frac{(\Delta-1)^{g/2}-1}{\Delta-2}, & g \text{ sodo}.
        \end{cases}
    \end{equation*}
        \end{trditev}
\end{frame}


\begin{frame}
\frametitle{Geometrija in posplošeni večkotniki}
\begin{definicija}
Geometrija ranga 2 \(\Gamma=(\PP, \LL, \II)\) je trojica 
nepraznih disjunktnih množic točk $\PP$, premic $\LL$, ter 
incidenčne relacije \(\II \subseteq \PP \times \LL\).
\end{definicija}
\begin{definicija}
$\Gamma=(\PP', \LL', \II')$ je podgeometrija geometrije $\Gamma=(\PP, \LL, \II)$, če velja
\(\PP' \subseteq \PP\), \(\LL' \subseteq \LL\) in
 \(\II' = \II \cap (\PP' \times \LL)\).
\end{definicija}
\end{frame}


\begin{frame}
    \frametitle{Geometrija in posplošeni večkotniki}
    \begin{definicija}
    Običajni $n$-kotnik je geometrija ranga 2, sestavljena iz 
    $n$ točk in $n$ premic, kjer se z incidenčno relacijo povežejo 
    v cikel dolžine $2n$.
    \end{definicija}
    \begin{definicija}
    Posplošeni $n$-kotnik je geometrija ranga 2 za katero velja,
    \begin{itemize}
        \item[\textbullet] da ne vsebuje $k$-kotnika z $2 \leq k < n$ (kot podgeometrijo),
        \item[\textbullet] poljubna $x, y \in \PP \cup \LL$ sta vsebovana v nekem 
        $n$-kotniku (kot podgeometrijo).
        \item[\textbullet] v geometriji obstaja vsaj en $(n+1)$-kotnik (kot podgeometrija).
    \end{itemize}
    \end{definicija}
\end{frame}

\begin{frame}
    \frametitle{Geometrija in posplošeni večkotniki}
    \begin{definicija}
    Naj bo \(\Gamma=(\PP,\LL,\II)\) geometrija ranga 2,
    njen incidenčni graf $G$ je dan z 
     množico vozlišč \(V(G) = \PP \cup \LL\) in množico 
    povezav \(E(G) = \{ uv \mid u \II v \}\).
    \end{definicija}
\end{frame}

\end{document}
